\section{Department priorities for future action}
\label{sec:priorities}

\subsection{Identify the department’s key issues relating to gender equality and establish key priorities for action over the next four years:} \instr{
\begin{itemize}
    \item Select up to five key priority areas where the department will strive for impact. Selected priorities should be justifiable and based on the quantitative and qualitative evidence presented in Section 2.
\item	Specific action(s) to support progress in priority areas should be identified. 
\end{itemize}
}





\subsubsection{Select up to five key priority areas where the department will strive for impact.
Selected priorities should be justifiable and based on the quantitative and qualitative
evidence presented in Section 2. Specific action(s) to support progress in priority areas should be identified.}




%\begin{itemize}[label=\textbf{\arabic*}]

\paragraph{Challenge}
The school has no formal structures nor has had discussions relating to a culture of gender equality and diversity. 

\begin{priority}{1}
 Develop the governance, structures, policies and procedures to deliver the AS action plan and imbed the AS Ireland Principles in the school.
\hfill{
    \hypersetup{hidelinks}
    \hyperlink{p1}{\gotoactionsymbol{}}
}
\end{priority}


A first action is to establish a gender balanced EDI Implementation Committee (EDIIC), and appoint a Chair from amongst the senior staff (Action~\ref{action:intro:edi}). The committee will have the authority to oversee the implementation of the action plan and develop a positive EDI culture in the school. 


\paragraph{Challenge}
Computer Science is a male-dominated discipline; Priorities \hyperlink{priority:2}{2}  and \hyperlink{priority:3}{3} will address the gender imbalance in CSIT staff (23\% F) and student (26\% F) cohorts.

\begin{priority}{2}
    Increase the ratio of female applicants for staff posts (Academic / Research) to improve gender balance.  
    \hfill{
    \hypersetup{hidelinks}
    \hyperlink{p2}{\gotoactionsymbol{}}
}
\end{priority}


%The main actions include improving the advertising materials to appeal to female applicants and to use female targeted initiatives.   

\begin{priority}{3}
    Increase the number of female students registered in UG and PG programmes. 
    \hfill{
    \hypersetup{hidelinks}
    \hyperlink{p3}{\gotoactionsymbol{}}
}
\end{priority}

Gender imbalance is more pronounced at undergraduate level, particularly in our flagship BScCS programme (14\%). The school will survey our UG students to understand what motivated them to study Computer Science (Action~\ref{action:survey:females}). It will enrich opportunities for female students to gain work experience, for example, through internships in the school (Action~\ref{action:enrich:opportunities}). Boosting the visibility of successful female researchers, alumni, and professionals will provide role models for new recruits and current students (Action~\ref{action:boost:visibility}).  

The actions associated with this priority area will provide information to improve our future recruitment strategies, deepening the pipeline for research and academic careers in Computer Science. 

\paragraph{Challenge}
Due to the lack of formal Professional Development Reviews (PDR) and mentoring in recent years, more support is required for staff career development and progression. 
There is a dearth of female academic (9\%) and research (17\%) staff in the school, with none at senior grades.  The school had no permanent female academic staff member until 2017.


\begin{priority}{4}
Supporting staff development to further career progression.
\hfill{
    \hypersetup{hidelinks}
    \hyperlink{p4}{\gotoactionsymbol{}}
}
\end{priority}


  
The key actions to deliver on this priority include restarting the Performance Development Review System to support staff career development (Actions~\ref{action:acad:pdrs}, \ref{action:acad:mentorresearchers}). Improved mentoring for interested staff, career training and research funding workshops will be conducted (Actions~\ref{action:acad:mentoring}, \ref{action:acad:supportfunding}).  

\paragraph{Challenge}
Staff would like to see more family-friendly policies, more diversity and inclusivity in staff and student cohorts.

\begin{priority}{5}
Promote a culture of inclusivity and belonging.
\hfill{
    \hypersetup{hidelinks}
    \hyperlink{p5}{\gotoactionsymbol{}}
}
\end{priority}

The key actions under this priority are both inward and outward focused.  The inward actions include supporting family / flexible leave and more integration of all research groups in school events (Actions~\ref{action:cult:famleave}, \ref{action:cult:researchinteraction}).  The outward facing actions concentrate on ensuring that schools from diverse backgrounds are a focus of attention (Action~\ref{action:cult:backgrounds}).




\subsection{Outline how the department’s gender equality priorities align with the institution’s Athena Swan action plan and, where relevant, broader EDI initiatives in the institution and/or department. This should include comment on:}
\instr{
\begin{itemize}
\item	key institutional actions that have, or will, support the department’s progress; 



\item	any gaps in institutional supports for achieving progress and impact in the department. 
\end{itemize}
}


\subsubsection{Comment on key institutional actions that have, or will, support the department’s progress}

The university bronze Athena Swan application (2019) has supporting actions which the school can leverage.\footnote{\url{https://www.ucc.ie/en/media/support/edi/athenaswan/AthenaSWANUCCActionPlan2019-2023.pdf}} 
Notably, the university aims to ‘Increase the ratio of female applicants for staff posts (Academic / Research) through the following actions. 
Action \hyperlink{act515}{5.1.5} will support our \hyperlink{priority:1}{Priority 1}:

\begin{UCCaction}{5.1.5}{}\hypertarget{act515}
Strengthening statements of commitment to EDI principles and practice incorporated into the President’s welcome information provided to applicant and a target of 40\% female applicants for academic post by 2023.
\hfill{
    \hypersetup{hidelinks}
    \href{https://www.ucc.ie/en/media/support/edi/athenaswan/AthenaSWANUCCActionPlan2019-2023.pdf}{\gotoactionsymbol{}}
}
\end{UCCaction}



The university commits to improving promotion pathways for senior academic staff (Actions \hyperlink{act513}{5.1.3} and \hyperlink{act519}{5.1.9}).  The university's success criterion and outcome is stated as: \textit{``Promotions across the University are at a minimum of 40\% female in all SL and Professor Scale 2 calls''}; this, in combination with a planned promotion pathway to Professor Scale 1,
is encouraging for CSIT academic staff (\hyperlink{priority:2}{Priority 2}). 

\begin{UCCaction}{5.1.3}{}\hypertarget{act513}
Establish a promotion pathway to Professor Scale 1, drawing on learnings for gender equality implemented in schemes for promotion to SL (Jan 2019) and Prof Scale 2 (Oct 2019) scheme calls (see Action 5.1.4).\\[.1in]
\end{UCCaction}

\begin{UCCaction}{5.1.4}{}\hypertarget{act514}
Submit applications for professorial appointments to the HEA Senior Academic Leadership Initiative (SALI).
\end{UCCaction}

\begin{UCCaction}{5.1.9}{}\hypertarget{act519}
Following the two promotion calls for Senior Lecturer and Professor (Scale 2) in 2019 in the University’s revised promotion schemes, review the number and success of female applicants to assess the gender impact of the revise teaching, research and service criteria and their weighting, and provision for and the provision for statutory leave in these schemes.
\end{UCCaction}


Action \hyperlink{act516}{5.1.6} in the university application addresses the need for recruitment and selection training for recruitment of research staff (Action~%2.2.1
\ref{action:acad:recruittraining}).  

\begin{UCCaction}{5.1.6}{}\hypertarget{act516}
Pilot action in SEFS to require Principal Investigators (PIs) participate in LEAD equality e-learning (including module on recruitment \& selection training) prior to submitting for approval a new funding application which involves recruitment of research staff.
\end{UCCaction}

Many of the actions under the university ‘Supporting and Advancing Careers’ will support the school’s ambitions in to support our staff develop their careers (\hyperlink{priority:4}{Priority 4}). 

The university's Action \hyperlink{act552}{5.5.2} will support the school to ensure that line managers know their responsibilities regarding family related leave and flexible working (\hyperlink{priority:5}{Priority 5}).  

\begin{UCCaction}{5.5.2}{}\hypertarget{act552}
Implement training for new and existing line managers to outline their roles, responsibilities regarding family related leave and flexible working.
\end{UCCaction}


\subsubsection{Comment on any gaps in institutional supports for achieving progress and impact in the department}

The university does not have the same gender imbalance in academic staff (45\% F), research staff (53\% F) and students (58\% F) as experienced in the school, therefore this is not addressed at university level.